% Generated draft: requires mathematical and source review.
\documentclass[11pt]{article}
\usepackage[margin=25mm]{geometry}
\usepackage{fontspec}
\setmainfont{Latin Modern Roman}
\usepackage{amsmath,amssymb,mathtools,unicode-math}
\setmathfont{Latin Modern Math}
\usepackage{enumitem,needspace,xurl,hyperref}
\setlength{\parindent}{0pt}
\setlength{\parskip}{5pt}
\setlength{\emergencystretch}{3em}
\newcommand{\N}{\mathbb{N}}
\newcommand{\Z}{\mathbb{Z}}
\newcommand{\Q}{\mathbb{Q}}
\newcommand{\R}{\mathbb{R}}
\newcommand{\C}{\mathbb{C}}
\newcommand{\F}{\mathbb{F}}
\newcommand{\A}{\mathbb{A}}
\newcommand{\PP}{\mathbb{P}}
\newcommand{\E}{\mathbb{E}}
\newcommand{\eps}{\varepsilon}
\newcommand{\dd}{\,\mathrm{d}}
\begin{document}
\section*{UnsolvedMath ID: 30003590}
\textbf{Source problem code:} catalogue entry
\subsection{Definitions and mathematical statement}
Let \(X\) be a projective variety, and let \(d\) denote the dimension of the subvarieties under consideration. Write
\[
N_d(X)=\bigl(Z_d(X)\otimes_{\Z}\R\bigr)/\!\equiv_{\mathrm{num}},
\]
where \(Z_d(X)\) is the group of algebraic \(d\)-cycles and \(\equiv_{\mathrm{num}}\) denotes numerical equivalence. Equip this finite-dimensional real vector space with its usual topology. The source does not specify the base field, the range of \(d\), or additional hypotheses on \(X\).

Let \(\mathcal A_d(X)\) denote the collection of \(d\)-dimensional subvarieties of \(X\) that are ample in the sense intended by the source. That sense is not defined in the supplied record and must be fixed before the question has a precise interpretation. Under the convention that a generated convex cone includes the origin, set
\[
C_d^{\mathrm{amp}}(X)=
\left\{\sum_{i=1}^{m}a_i[Y_i]_{\mathrm{num}}:
 m\in\Z_{\geq 0},\quad a_i\in\R_{\geq 0},\quad
 Y_i\in\mathcal A_d(X)\right\},
\]
where the sum for \(m=0\) is \(0\). No closure is taken.

\textbf{Source question.} For every projective variety \(X\) and every dimension \(d\) covered by the intended hypotheses, is the convex cone generated by the numerical classes of ample \(d\)-dimensional subvarieties open in \(N_d(X)\)?

\textbf{Unresolved qualifications.} With the displayed cone convention, the question asks whether \(C_d^{\mathrm{amp}}(X)\) itself is open, including at the origin. The supplied wording does not establish whether the intended object instead excludes the origin, for example \(C_d^{\mathrm{amp}}(X)\setminus\{0\}\). These are different formulations and must not be identified silently. Likewise, openness in the ambient space \(N_d(X)\) is distinct from relative openness in the linear span of the cone. The wording names \(N_d(X)\), but supplies no explicit topology convention beyond this. The definition of ample subvariety, base-field assumptions, admissible dimensions, and any boundary cases such as \(d=0\) or \(d=\dim X\) require clarification. This is a statement-only draft, with these qualifications unresolved.
\subsection{Short English statement}
The supplied question asks whether numerical classes of ample subvarieties of a fixed dimension generate an open cone in the real numerical cycle space of a projective variety. A precise formulation requires the intended definition of ample subvariety, the hypotheses on the variety and dimension, and a convention about inclusion of the origin. The record labels the problem open; that status is supplied curation and has not been independently verified.
\subsection{Sources}
\begin{itemize}
\item Supplied local UnsolvedMath registry record, permanent ID \(30003590\), exact source problem code \texttt{catalogue entry}, title: \emph{Openness of Cones Generated by Ample Subvarieties}. Collection: Oberwolfach Reports Open Problems. No individual proposer is identified; the supplied proposed year is 2017.
\item Supplied report citation: \emph{Mini-Workshop: Positivity in Higher-dimensional Geometry: Higher-codimensional Cycles and Newton--Okounkov Bodies} (2018), Oberwolfach Reports, pp.~2631--2657, DOI \texttt{10.4171/owr/2017/43}. \url{https://doi.org/10.4171/owr/2017/43}. The background describes Oberwolfach Reports as documenting mathematical workshops and their open-problem discussions. The cited publication year 2018 and the supplied proposed year 2017 are retained without resolving their relationship.
\item Supplied literature source: Fulger--Lehmann, positive cones of dual cycle classes. \url{https://content.algebraicgeometry.nl/2017-1/2017-1-001.pdf}.
\item Supplied literature source: Fulger--Kuronya--Lehmann, OWR report. \url{https://doi.org/10.4171/OWR/2017/43}.
\item Literature assessment checked on \texttt{2026-08-21}, according to the supplied record: status open. The supplied assessment says that higher-codimension positivity results give meaningful full-dimensional positive cones generated by ample-subvariety classes, but that no general theorem identifies the cone generated by those classes with an open cone in the numerical cycle space. This is supplied dated curation derived from the OWR record, not an independently verified literature assessment; its status and cited claims require independent verification.
\item The registry labels the statement \texttt{corrected\_verified} and says that notation, typography, and sentence structure were repaired while preserving mathematical content. This is a supplied editorial qualification, not independent mathematical verification of this draft or of the cited report.
\end{itemize}
\end{document}
